\section{Exhaustion of the alternating Heisenberg quotient and the faithful partner}
\label{sec:heisenberg-alternating-quotient-exclusion}

\paragraph{Status and provenance.}
D-TD43, reserved question Q2ey, combines the separate primitive,
block-permutation and block-diagonal carrier theorems. Root proposed
the exhaustive combination and its main maximal-layer application;
\texttt{/root/g34\_eightfolds} supplies the complete interface check
below. The combination introduces no new character estimate,
classification input or numerical computation. Complete author, root and
independent whole-proof and dependency-interface readings pass. Exact
source scopes, shared provenance and review histories remain recorded.

\subsection*{D-TD43. No large alternating quotient of an irreducible label image}

Let an actual connected finite $J$-cover of the sphere have a finite
complex coefficient representation on a stable tensor product
$V=W\otimes R$, with the entire external factor $R$ retained. Suppose
the actual projective image on $W$ has the following data:
\[
 \begin{gathered}
 e=\dim V,\qquad
 \sum_i\operatorname{rank}(1-V(g_i))=2e+6,\\
 P=\operatorname{im}(J\to\operatorname{PGL}(W)),\qquad
 Q\triangleleft P,\qquad Q\simeq\mathbf F_p^{2m},\qquad m\ge1,\\
 H=P/Q\hookrightarrow\operatorname{Sp}(Q),\qquad
 \psi:H\twoheadrightarrow T\in\{A_n,S_n\}.
 \end{gathered}
 \tag{HAQ.1}
\]
Here the projective Heisenberg action of $Q$ on $W$ is nondegenerate
and irreducible, so $\dim W=p^m$. The actual defining
$\mathbf F_pH$-module $Q$ is irreducible. The maps $J\to P\to H\to T$
are actual surjections; the displayed embedding of $H$ is its actual
faithful symplectic action. An alternating subgroup or composition
factor is not substituted for the specified surjection $\psi$.

\begin{theorem}[D-TD43]
The actual source (HAQ.1) is impossible whenever
\[
                              n\ge2^{30}.
 \tag{HAQ.2}
\]
There is no assumption of absolute irreducibility over $\mathbf F_p$,
absolute primitivity or tensor indecomposability of its defining
module. There is no bound or support-floor hypothesis on the entire
complex partner $R$.
\end{theorem}

\begin{proof}
Use the intrinsic commuting field, not a chosen field of definition
for an absolute constituent:
\[
 F=\operatorname{End}_{\mathbf F_pH}(Q)=\mathbf F_q,\qquad
 q=p^f,\qquad a=\dim_F Q,\qquad af=2m,\qquad
 M=Q\otimes_F\overline{\mathbf F}_p.
 \tag{HAQ.3}
\]
The complete finite-algebra argument of D-TD34 makes $F$ a finite
field, makes every element of $H$ $F$-linear, and proves that $M$
is irreducible. Its uniform field bound also applies to (HAQ.1),
without any primitivity premise. Thus this is exactly the common
actual setup of D-TD38, D-TD40(b) and D-TD42. The original sphere
cover, its inertia operators and its complete complex tensor partner
are the same for each of these three applications.

If $M$ is primitive, D-TD38 excludes (HAQ.1) at (HAQ.2). That
theorem explicitly permits tensor decomposition of the absolute
module and does not require it to descend as an absolutely
irreducible prime-field module.

Otherwise choose a finest nontrivial $H$-stable direct-sum system
in $M$. Such a system exists: there are at most $a$ nonzero summands,
and a proper refinement increases their number. Write
\[
 M=\bigoplus_{j=1}^t U_j,\qquad t\ge2,\qquad
 K_{\mathrm{bl}}=
 \ker\bigl(H\to\operatorname{Sym}\{U_1,\ldots,U_t\}\bigr).
 \tag{HAQ.4}
\]
Irreducibility makes $H$ transitive on these blocks and their
dimensions equal. The blocks are allowed to be defined only over
the algebraic closure. The subgroup $K_{\mathrm{bl}}$ is normal
in the actual group $H$, so $\psi(K_{\mathrm{bl}})$ is normal in
$T$. For $n\ge5$, every nontrivial normal subgroup of $A_n$ or
$S_n$ contains $A_n$. Therefore exactly the following dichotomy
is exhaustive:
\[
 \psi(K_{\mathrm{bl}})=1
 \qquad\hbox{or}\qquad
 \psi(K_{\mathrm{bl}})\supseteq A_n.
 \tag{HAQ.5}
\]
In the first case $\psi$ factors through the actual transitive
permutation image on these $t$ blocks. D-TD40(b) excludes the
source, since $n\ge2^{30}>10^6$. Its proof permits every
transitive block-permutation image; no primitivity of that image
is being assumed here. In the second case the chosen finest
system satisfies the exact hypothesis of D-TD42, which also
excludes the source at (HAQ.2). That theorem needs no transitivity
of $\ker\psi$ on the blocks.

These cases cover the absolute irreducible module $M$. All three
input theorems retain the original finite periods and actual quotient
covers; none assigns a sphere source to a block stabilizer, a terminal
tensor coordinate or an abstract copy of $T$. Their faithful-image
passage also permits the coefficient representation in (HAQ.1) to
have a kernel: a tensor operator acting identically on $V$ acts
projectively trivially on $W$, so that kernel kills $P$, $H$ and
$T$. Passing to its actual connected quotient cover preserves the
moving-rank sum and all required quotient maps. Thus no hidden
faithfulness or field-descent premise remains in this combination.
\end{proof}

\subsection*{The maximal standard layer has a nonfaithful partner}

\begin{corollary}[D-TD43, maximal-layer consequence]
Use the actual finite primitive rank-nine coefficient-source setting
of HD-ML02 and one fixed complete D-TD12 decomposition. Suppose
its maximal alternating natural degree is $n\ge N_*$, with the
integer $N_*$ of (ML.1). For any selected maximal layer write its
actual finite linear factorization as
\[
 V=(\operatorname{Std}_n\circ x)\otimes U,\qquad
 x:J\twoheadrightarrow T\in\{A_n,S_n\},\qquad
 K_0=\ker U.
 \tag{HAQ.6}
\]
Then $U$ is not faithful. Moreover $x|_{K_0}$ is injective and
\[
                    A_n\subseteq x(K_0)\subseteq T.
 \tag{HAQ.7}
\]
In particular $K_0$ is transitive on the $n$ natural letters.
The conclusion concerns the actual linear partner of (HAQ.6),
including all earlier and later pieces of the fixed recursion.
\end{corollary}

\begin{proof}
HD-ML02 says that the selected maximal extraction layer is a single
ordinary standard factor; its rational standard lift gives the actual
linear factorization (HAQ.6), as in HD-SS01. The complete orbit bound
D-TD16 applies to the same chosen decomposition.

The comparison of thresholds is elementary and explicit. From (ML.1),
both $F_*\ge241$ and $100(1+M_*)\ge241$. Consequently
\[
 N_*=[100(1+M_*)F_*^{240}]^2
       \ge241^{482}>2^{482}>2^{30}.
 \tag{HAQ.8}
\]
If $U$ were faithful, D-TD19 would supply an original coarse
Heisenberg coordinate in $U$ with an actual surjective chain
$J\to P_j\to H_j\to T$. Its defining prime-field module is
irreducible, its whole Heisenberg space is stable, and the complement
of that space in $V$ remains an entire tensor partner. The original
sphere cover and rank sum remain unchanged. Thus this coordinate
satisfies (HAQ.1), and (HAQ.8) puts it in the impossible range of
D-TD43. Hence $K_0\ne1$.

If $k\in K_0$ also satisfies $x(k)=1$, then both factors in
(HAQ.6) act identically on $k$. The coefficient image $J$ is
faithful on $V$, so $k=1$. This proves injectivity of $x|_{K_0}$.
Its image is therefore a nontrivial normal subgroup of $T$,
which proves (HAQ.7) and transitivity. Neither this argument nor
D-TD19 assumes independence of distinct quotient actions or lifts
later extraction groups to subgroups of the original $J$.
\end{proof}

\paragraph{Exact remaining minimal-source gate.}
Assume in addition the exact imaginary-quadratic $K$-model,
signature $(3,3)$, actual period rank nine and restricted Jacobian
polarization of HD-SS06. Put $b=\dim U\ge2$, let
$J_1=x^{-1}((S_n)_1)$, and retain the actual cohomology multiplicity
\[
 \mu_U=\dim_{\mathbf C}
          \operatorname{Hom}_J(U,H^1(C,\mathbf C)).
 \tag{HAQ.9}
\]
Suppose the source has smallest coefficient degree among the actual
exact-$K$, rank-nine sources with its natural signed polarization
determinant class, exactly as in HD-SS11. Then every surviving
configuration in (HAQ.6) satisfies
\[
 \begin{gathered}
 K_0\ne1,\quad x(K_0)\supseteq A_n,\quad
 U(J_1)=U(J),\quad F_U=F_{U|_{J_1}}=K,\\
                         \mu_U>0,\qquad
                 n\le72(b+1),\qquad b\ge n/72-1.
 \end{gathered}
 \tag{HAQ.10}
\]
Here $K_0$ is a group kernel and $K$ is a coefficient field.

Indeed (HAQ.7) gives a normal subgroup mapping onto $A_n$ and
acting trivially on $U$. The stronger actual-subgroup interface of
HD-SS01 gives the image and field equalities in (HAQ.10), including
models over $K$ and Schur index one. If $\mu_U=0$, the post-vanishing
part of HD-SS06 gives an actual rational-base source of coefficient
degree $b$ with the same exact field, signature and period rank.
Its specified multiplicity-one Rosati coefficient-line correspondence
changes the rank-six Hermitian polarization only by a positive
rational scalar, whose sixth power is a norm. The natural signed
determinant class is unchanged. That post-vanishing argument requires
transitivity and $\mu_U=0$, not its earlier sufficient numerical
condition $n>168(b+1)$. Since $b<(n-1)b$, minimality is contradicted.
Thus $\mu_U>0$. HD-SS11 now yields $n\le72(b+1)$, giving the final
inequality of (HAQ.10).

The positive $U$-cohomology summand is therefore retained. No upper
bound on $b$, $n$, the remaining inducing image or an original
overgroup follows from (HAQ.10). The theorem excludes the faithful
partner branch and exhausts the specified actual alternating
Heisenberg quotient; it does not assert a finite complete catalogue,
a new detecting cycle, or nonexistence of the kernel-transitive
large-partner branch. Without the stated minimality, no vanishing or
positive-cohomology conclusion is substituted for the exact HD-SS02
direct-sum identity.
